\documentclass[planificacion2024.tex]{subfiles}
\begin{document}

\section*{\centering PROGRAMA DE ESTUDIOS}

\vspace{2cm}
\begin{tabular}{r l}
\textbf{ESPACIO CURRICULAR:}& ELECTROTECNIA II\\
\textbf{CURSO :}& 4° 1°\\
\textbf{PROFESOR/A:}    & FERREYRA GUSTAVO DAVID \\
\textbf{TITULO DEL EGRESADO:}&Técnico en Electricidad
\end{tabular}

\vspace{1cm}

Electrotecnia II se fundamenta en la necesidad de profundizar 
los saberes y aprendizajes relacionados con la asociación y comportamiento técnico de los capacitores desde un punto de vista de carga eléctrica y capacidad pertinente. 
Profundiza los saberes relacionados con la Inducción Electromagnética, que serán retomados en años posteriores en la comprensión del funcionamiento de máquinas eléctricas con características electrostáticas y electrodinámicas, donde el magnetismo y la electricidad se hacen presentes dando origen a transformaciones energéticas. 

Además, se desarrollarán conocimientos eléctricos vinculados con los principios fundamentales de origen y aplicación de la energía alterna. 

Se generan propuestas que impliquen resolver esquemas 
eléctricos y situaciones problemáticas energéticas, relacionadas con temáticas de electromagnetismo, de empleo de energía alterna y análisis crítico del comportamiento de capacitores conectados con fuentes de alimentación eléctrica en un marco resistivo-capacitivo. 
\begin{enumerate}
    \item CIRCUITOS RESISTIVO- CAPACITIVOS 
    \begin{itemize}
        \item Conocer los procedimientos de cálculo en un circuito con características capacitivas transitorias.
        \begin{itemize}
            \item Caracterización de la obtención analítica de la constante de tiempo, en circuitos resistivo-capacitivos, con corriente eléctrica continua. 
            \item Interpretación de fenómenos de carga y descarga en un circuito en un esquema de aplicación resistivo-capacitivo. 
            \item Análisis comparativo de las características de las curvas de carga-descarga energética, utilizando software específico. 
            \item Comprobación y síntesis de cálculos eléctricos analíticos, relacionados con componentes capacitivos energizados, utilizando software específico. 
        \end{itemize}
        
        \item Interpretar las teorías electrodinámicas eléctricas en una máquina rotativa.
    \end{itemize}
    \item INDUCCIÓN ELECTROMAGNÉTICA 
    \begin{itemize}
        \item Analizar de modo crítico las 
leyes electromagnéticas 
relacionadas con la interacción entre circuitos 
con alimentación eléctrica. 
    \begin{itemize}
        \item Conceptualización de las leyes de Faraday-Lenz 
relacionadas con la fuerza electromotriz en bobinas inductivas. 
        \item Reconocimiento de las características ferromagnéticas de 
los circuitos de aplicación según Hopkinson. 
\item Interpretación de las características magnéticas 
desarrolladas por Biot-Savart. 
\item Distinción entre las conclusiones electromagnéticas 
enunciadas por Gauss, Maxwell, Tesla y Weber, respecto 
del comportamiento de los circuitos frente a la influencia 
de electricidad en esquemas ferromagnetizables. 
    \end{itemize}
        \item Conocer los principios de 
inducción y autoinducción 
electromagnética desde un 
punto de vista de 
interacción energética. 
\begin{itemize}
    \item Caracterización de los fenómenos de inducción 
magnética e inducción mutua en circuitos con 
alimentación eléctrica. 
\item Caracterización de la reluctancia magnética y la fuerza 
magneto-motriz desde una perspectiva de interacción
entre ambos fenómenos. 
\item Interpretación de la autoinducción como fenómeno 
electromagnético propio, vinculado con las características 
intrínsecas de una bobina energizada. 
\item Formalización de las reglas prácticas de uso anatómico. 
\item Aplicación de los sistemas de unidades C.G.S. y/o
Internacional desde una óptica magnética, en la 
resolución de circuitos de modelaje electromagnético. 
\item Aplicación de las leyes electromagnéticas a los circuitos 
resistivos-capacitivos e inductivos. 
\end{itemize}
    
    \end{itemize}
    \item PRINCIPIOS FUNDAMENTALES DE LA ENERGÍA ALTERNA
    \begin{itemize}
        \item Conocer la naturaleza de 
generación de energía 
eléctrica alterna. 
\begin{itemize}
    \item  Reconocimiento del origen técnico-eléctrico de energía 
alterna. 
\item Interpretación gráfica de la onda senoidal, desde un punto 
de vista eléctrico con energía alterna presente. 
\item Formalización de la energía alterna. 
\end{itemize}
        \item Comprender el uso de la 
energía alterna en el 
mundo energético actual.
\begin{itemize}
    \item Argumentación de las características intrínsecas de la 
energía eléctrica alterna. 
\item Formulación de los valores típicos representativos de 
origen energético alterno senoidal. 
\item Conceptualización de las fórmulas técnicas eléctricas, que 
son de aplicación en el estudio de la energía alterna.
\end{itemize}
    \end{itemize}
\end{enumerate}

\end{document}